\documentclass[10pt,a4paper]{amsart}
\usepackage[margin=19mm]{geometry}
\usepackage{amsmath,amssymb,mathtools}
\usepackage[scr=rsfs,cal=euler]{mathalfa}
\usepackage{xfrac}
\usepackage[unicode,hidelinks,bookmarks=false]{hyperref}
\newcommand{\ie}{i.e.\ }
\newcommand{\cf}{cf.\ }
\newcommand{\resp}{respectively\ }
\newcommand{\Subsection}{\subsection}
\providecommand{\nobreakdash}{\nobreak\hbox{-}}
\setlength{\emergencystretch}{2em}
\multlinegap0pt
\usepackage{bm}
\usepackage{bbm}
\usepackage{dsfont}
\usepackage{xspace}
\usepackage{cancel}
\usepackage{mathtools}
\usepackage{amsthm}
\makeatletter
\@ifundefined{theorem}{\newtheorem{theorem}{Theorem}}{}
\@ifundefined{lemma}{\newtheorem{lemma}[theorem]{Lemma}}{}
\@ifundefined{notation}{\newtheorem{notation}[theorem]{Notation}}{}
\makeatother
\makeatletter
\expandafter\ifx\csname example\endcsname\relax
\newenvironment{example}
{\pushQED{\qed}\renewcommand{\qedsymbol}{$\diamond$}\examplex}
\fi
\makeatother
\title[An edge-bicolored graph approach to the Ising model on rando]{An edge-bicolored graph approach to the Ising model on random regular graphs}
\author{Michael Borinsky, Shiyue Ren, Maximilian Wiesmann}
\date{Source: arXiv:2607.07867v1 (CC BY 4.0)}
\begin{document}
\maketitle

\noindent\emph{Source and licence.} An edge-bicolored graph approach to the Ising model on random regular graphs. Version arXiv:2607.07867v1 (CC BY 4.0), distributed under the Creative Commons Attribution 4.0 International licence, as recorded in the arXiv licence metadata for this exact version and shown on the abstract page https://arxiv.org/abs/2607.07867v1. Full rights notice: licenses/arxiv-2607.07867-CC-BY-4.0.txt.\par\smallskip

\noindent\textit{Editorial scope.} Counted unit: the Lemma whose source label is \texttt{lem:maximum} in the article (source0.tex lines 460--479, source label \texttt{lem:maximum}), delivered together with the article's own complete original proof of it (source0.tex lines 481--528, one \texttt{proof} environment of 48 lines). No mathematics is written, repaired or re-derived: statement and proof are the article's own text line for line. Required context retained uncounted: not:variable\_changes (lines 252--255); lmm:positiveVk (lines 431--435); eq:crit (lines 449--452); eq:gdef (lines 454--457). Every definition, display and label of those lines is retained unchanged. Omitted from this package and not redistributed: 1--251, 256--430, 436--448, 453--453, 458--459, 480--480, 529--868. Nothing inside a retained excerpt is omitted or altered: every retained body is delivered exactly as the article's lines are written, so no omission receipt of any kind accompanies this package. Citations: the retained excerpts carry no citation command at all, so no bibliography entry of the article is transcribed and no reference is redistributed. Reference adaptation: none was needed and none was made; every reference inside the delivered unit resolves inside it, so no unresolved reference or citation is produced. Printed slips kept literal: no printed word, symbol, hypothesis or number of the counted statement or of its proof was altered. Layout and class: the article's own class is \texttt{article} with packages including hyperref, geometry, inputenc, fontenc, lmodern, microtype, graphicx, amsmath, amssymb, amsthm, xcolor, biblatex; the wrapper is the lane's pinned \texttt{amsart} editorial wrapper at 10pt on A4, which restates the theorem-like environments the retained text needs and adds the article's own macro definitions that the retained text needs (none), with extra wrapper packages (bm, bbm, dsfont, xspace, cancel, mathtools). The delivered numbering is the unit's own. Assets: no retained span contains a figure environment, a graphics-inclusion command, a PDF-page inclusion, a table or a TikZ picture; no image file is copied into this package, the package declares no asset dependency and no image file is redistributed with it. Licence: version arXiv:2607.07867v1 of this article is distributed under the Creative Commons Attribution 4.0 International licence, as recorded in the arXiv licence metadata for this exact version and shown on the abstract page https://arxiv.org/abs/2607.07867v1. A full rights notice is delivered at licenses/arxiv-2607.07867-CC-BY-4.0.txt. Characters: every retained excerpt is pure ASCII, so the delivered engine, XeLaTeX with its default Latin Modern text font, reports no missing character.
\begin{notation}
\label{not:variable_changes}
We define $\kappa = \tanh(\beta J)$, $\tau = \tanh(\beta h)$, $\rho = \cosh(\beta J)$ and $\eta = \cosh(\beta h)$. Note that high temperature $T > T_c$ means $\kappa < \kappa_c = \tanh\big(\tfrac{1}{2} \log \big(\tfrac{k}{k-2}\big)\big) = \tfrac{1}{k-1}$.
\end{notation}

\begin{lemma}
\label{lmm:positiveVk}
If $\kappa, \tau \geq 0$, then we have $|V_k(x,y)| \leq V_k(|x|, |y|)$ for all $(x,y) \in S^1$.
Moreover, if $\kappa \geq 0$ and $\tau \leq 0$, then $|V_k(x,y)|  \leq V_k(|x|,-|y|)$ for all $(x,y) \in S^1$. \qed
\end{lemma}

\begin{align}
  \label{eq:crit}
  g(\theta^\star) = \frac{1-\tau}{1+\tau}\, ,
\end{align}

\begin{align}
  \label{eq:gdef}
  g(\theta) = \frac{\sin(\phi-\theta)}{\sin(\phi+\theta)} \left( \frac{\cos(\phi-\theta)}{\cos(\phi+\theta)} \right)^{k-1}\,,
\end{align}

\begin{lemma}\mbox{}
  \label{lem:maximum}
  \begin{enumerate}
    \item If $h > 0$, then $|V_k(\theta)|$ has a unique global maximum in the interval $[0,\pi/2]$ at $\theta= \theta^\star$.
    \item If $h < 0$, then $|V_k(\theta)|$ has a unique global maximum in the interval $[-\pi/2,0]$ at $\theta=-\theta^\star$.
  \end{enumerate}

  Furthermore, let $\beta_c^{\mathrm{tree}} = \frac{1}{2J}\log(\frac{k}{k-2})$.  If $h =0$ and
  \begin{enumerate}
    \item[3.] $\beta > \beta_c^{\mathrm{tree}}$, then $|V_k(\theta)|$ has a unique global maximum in the interval $[0,\pi/2]$ at $\theta = \theta^\star$.
    \item[4.] $\beta < \beta_c^{\mathrm{tree}}$, then $|V_k(\theta)|$ has a unique global maximum in the interval $[0,\pi/2]$ at $\theta = 0$.
  \end{enumerate}

  Here, $\theta^\star$ is the unique solution of 
  \begin{align}
    g(\theta^\star) = \frac{1-|\tau|}{1+|\tau|}\, ,
  \end{align}
  in the interval $(0,\pi/2-\phi)$ with $g$ defined in eq.~\eqref{eq:gdef}.
  (Recall that $\tau$ is a function of $h$ by Notation~\ref{not:variable_changes}.)
\end{lemma}

\begin{proof}
  The derivative of $g$ fulfills
  \begin{align*}
    g'(\theta) = \sin(2\phi)\frac{\cos(\phi-\theta)^{k-2}}{2\cos(\phi+\theta)^{k}\sin^2(\phi+\theta)} \bigl( (k-2) \cos(2\theta) - k \cos(2\phi) \bigr)\,,
  \end{align*}
  and the equation
  \begin{equation}
    \label{eq:theta_0}
    \cos(2\theta_0) = \frac{k}{k-2}\cos(2\phi)\, ,
  \end{equation}
  has a unique solution $\theta_0 \in [0,\pi/2]$ if and only if 
  $\cos(2\phi) \leq \frac{k-2}{k}$. If such a unique $\theta_0$ exists, then $g'(\theta_0)=0$. This condition is equivalent to $\kappa \geq \frac{1}{k-1}$ as $\cos(2\phi) = \frac{1-\kappa}{1+\kappa}$ 
  and also to $\beta \geq \frac{1}{J} \tanh^{-1}\left( \frac{1}{k-1} \right) = \beta_c^\mathrm{tree}$.

  If there is such a $\theta_0$, then $0<\theta_0 < \pi/4 < \theta_s$
  and $g$ is monotonically increasing in the interval $[0,\theta_0]$ and
  monotonically decreasing in the interval $[\theta_0,\theta_s)$. 
  On the other hand, we find that  if $\beta < \beta_c^\mathrm{tree}$, then
  $g$ is monotonically decreasing on $[0,\theta_s)$.

  On $(\theta_s,\pi/2]$, the derivative of $g$ is non-vanishing. Because $g$ has a pole at $\theta_s$, 
  its magnitude $|g(\theta)|$ must be decreasing on $(\theta_s,\pi/2]$.
  Because $|g(\pi/2)|=1$, it follows that $|g(\theta)| >1$ for $\theta \in (\theta_s,\pi/2)$.
  Furthermore, note that $g(0) =1$ and $\lim_{\theta \rightarrow \theta_s^-} g(\theta) = -\infty$.

  If $h > 0$, and therefore $0<\frac{1-\tau}{1+\tau} < 1$, \eqref{eq:crit} has a unique solution at $\theta^\star \in (0,\theta_s)$,
  as either $g$ is monotonically decreasing in that interval or $g$ is increasing from $g(0)=1$ to some maximal value $g(\theta_0) >1$ before it 
  decreases to $-\infty$ in the interval $[\theta_0,\theta_s)$. The case $h<0$ follows by the symmetry $(\tau,\theta) \leftrightarrow (-\tau,-\theta)$ of the potential $V_k$.
  There are no other critical points in the interval $[0,\pi/2]$. So, by Lemma~\ref{lmm:positiveVk}, the critical point of $V_k$ at $\theta^\star$ is a global maximum.

  This proves the $h \neq 0$ case of the statement.

  If $h=0$, \eqref{eq:crit} reduces to $g(\theta^\star) = 1$. 
  If, moreover, $\beta < \beta_c^\mathrm{tree}$, then \eqref{eq:crit} 
  has one or two solutions on the interval $[0,\pi/2]$ as $g(0) = 1$ and $g(\pi/2) = (-1)^k$.
  The discussion above implies that $g(\theta) \neq 1$ for all $\theta \in (0,\pi/2)$.
  It is easy to verify that $|V_k(0)| > |V_k(\pi/2)|$ as $\phi \in (0,\pi/4)$. Hence, only $\theta = 0$ is a
  global maximum of $V_k$ by Lemma~\ref{lmm:positiveVk}.

  If $\beta > \beta_c^\mathrm{tree}$, there is an additional candidate 
  for a global maximum at $\theta^\star \in (\theta_0,\theta_s)$, where \eqref{eq:crit} has a nontrivial solution.
  Computing the second derivative of $V_k$ and evaluating at $\theta=0$ gives
  \begin{equation}
      \label{eq:2nd_der_at_0}
      \frac{\partial^2}{\partial \theta^2} V_k(0) = \frac{2(1+\kappa)^{k/2}}{2 k!} k \cos^k(\phi) \left( (k-1) \tan^2(\phi) - 1 \right).
  \end{equation}
  Since $\beta > \beta_c^{\text{tree}} = \frac{1}{J}\mathrm{arctanh} \big(\tfrac{1}{k-1}\big)$, we have $\tan^2(\phi) = \kappa > \frac{1}{k-1}$. Moreover, $\cos^k(\phi) > 0$ as $0<\phi<\pi/4$, hence \eqref{eq:2nd_der_at_0} is positive and $\theta=0$ is a local minimum of $V_k$. Since $\theta=0$ and $\theta=\theta^\star$ are the only critical points, $\theta^\star$ is a local and hence global maximum.
\end{proof}

\end{document}
